% chapters/requirements/rf_notifications.tex
% RF-13: Notificaciones
% ============================================================================

\item \textbf{Notificaciones.} \label{rf:notifications}
El sistema de notificaciones opera en dos canales: notificaciones \textit{in-app}
(almacenadas en la base de datos y consultables por el usuario mediante \dfn{polling}) y
correo electrónico (espejo opcional de la notificación \textit{in-app}, activable por cada
usuario). El envío de correos se procesa de forma asíncrona mediante la cola de tareas.
Las notificaciones se archivan en los cambios de curso. Existen además correos directos
(bienvenida, restablecimiento de contraseña) que no generan notificación \textit{in-app}.

\begin{functional}

    % RF-13.1 ---------------------------------------------------------------
    \item \textbf{Creación de notificación \textit{in-app}.}
    \label{rf:notifications:create}
    El sistema debe crear notificaciones internas asociadas a eventos del dominio.

    \textbf{Entrada:} evento interno del sistema (publicación de notas, apertura de
    revisión, asignación de corrección, resultado de revisión, etc.).

    \textbf{Proceso:} el sistema crea un registro de notificación en la base de datos con el
    tipo de evento, el mensaje, la fecha, el usuario destinatario y el estado de lectura (no
    leída por defecto). El contenido del mensaje se genera a partir del tipo de evento y los
    datos del contexto.

    \textbf{Salida:} notificación almacenada y disponible para consulta por el usuario
    destinatario.

    % RF-13.2 ---------------------------------------------------------------
    \item \textbf{Envío de correo electrónico espejo.}
    \label{rf:notifications:email_mirror}
    El sistema debe enviar un correo electrónico espejo para cada notificación
    \textit{in-app}, si el usuario destinatario tiene habilitadas las notificaciones por
    correo.

    \textbf{Entrada:} evento interno disparado tras la creación de una notificación
    \textit{in-app}.

    \textbf{Proceso:} el sistema verifica la preferencia de notificación del usuario. Si el
    correo está habilitado, encola una tarea asíncrona en la cola de notificaciones para
    enviar el correo con la plantilla correspondiente al tipo de evento. El fallo en el
    envío del correo no afecta a la notificación \textit{in-app}.

    \textbf{Salida:} correo electrónico enviado al usuario, o registro de error en el
    \textit{log} de aplicación si el envío falla.

    % RF-13.3 ---------------------------------------------------------------
    \item \textbf{Marcado de notificación como leída.}
    \label{rf:notifications:read}
    El sistema debe permitir al usuario marcar sus notificaciones como leídas.

    \textbf{Entrada:} petición PATCH al recurso de la notificación (individual o masiva) con
    el nuevo estado de lectura.

    \textbf{Proceso:} el sistema actualiza el estado de lectura de la notificación o
    notificaciones indicadas. Se permite el marcado masivo (marcar todas como leídas).

    \textbf{Salida:} respuesta con el estado de lectura actualizado.

    % RF-13.4 ---------------------------------------------------------------
    \item \textbf{Listado de notificaciones.} \label{rf:notifications:list}
    El sistema debe permitir al usuario consultar sus notificaciones con paginación.

    \textbf{Entrada:} petición GET al recurso de notificaciones propias con parámetros
    opcionales de filtrado (estado de lectura, tipo de evento) y paginación.

    \textbf{Proceso:} el sistema recupera las notificaciones del usuario ordenadas por fecha
    descendente, aplicando los filtros indicados. Las notificaciones archivadas no aparecen
    en el listado por defecto.

    \textbf{Salida:} respuesta con la lista paginada de notificaciones.

    % RF-13.5 ---------------------------------------------------------------
    \item \textbf{Archivado de notificaciones en cambio de curso.}
    \label{rf:notifications:archive}
    El sistema debe archivar las notificaciones del curso anterior al ejecutar la transición
    de curso.

    \textbf{Entrada:} evento interno disparado por la transición de curso
    (\cref{rf:courses:transition}).

    \textbf{Proceso:} el sistema marca todas las notificaciones asociadas al curso archivado
    como archivadas. Las notificaciones archivadas dejan de aparecer en el listado por
    defecto pero siguen accesibles para los gestores.

    \textbf{Salida:} notificaciones archivadas.

    % RF-13.6 ---------------------------------------------------------------
    \item \textbf{Envío de correo directo.} \label{rf:notifications:direct_email}
    El sistema debe enviar correos electrónicos directos que no generan notificación
    \textit{in-app}.

    \textbf{Entrada:} evento interno de bienvenida con contraseña (usuario nuevo) o de
    restablecimiento de contraseña.

    \textbf{Proceso:} el sistema encola una tarea asíncrona para enviar el correo con la
    plantilla correspondiente. Estos correos se envían siempre, independientemente de las
    preferencias de notificación del usuario.

    \textbf{Salida:} correo electrónico enviado.

    % RF-13.7 ---------------------------------------------------------------
    \item \textbf{Plantillas de correo electrónico y notificación.}
    \label{rf:notifications:templates}
    El sistema debe definir y gestionar plantillas de correo para los distintos tipos de
    notificaciones.

    \textbf{Entrada:} no aplica directamente; este requisito define las plantillas que el
    sistema debe implementar.

    \textbf{Proceso:} el sistema debe implementar las siguientes plantillas con soporte de
    internacionalización: bienvenida con contraseña, restablecimiento de contraseña,
    publicación de notas, apertura de revisión, solicitudes de revisión agrupadas (para
    corrector), asignación de corrección (para corrector), resultado de revisión (para
    alumno) y recordatorio de borrado automático (para gestor). Las plantillas reciben los
    datos del contexto como variables.

    \textbf{Salida:} correos generados con el contenido y formato adecuados.

\end{functional}
